\documentclass[11pt]{article}
\usepackage{amsmath,amssymb,amsthm}
\usepackage[margin=1.2in]{geometry}
\newtheorem{theorem}{Theorem}
\newtheorem{lemma}[theorem]{Lemma}
\newtheorem{proposition}[theorem]{Proposition}
\newtheorem{corollary}[theorem]{Corollary}
\newtheorem{remark}[theorem]{Remark}
\newcommand{\e}[1]{e\!\left(#1\right)}
\newcommand{\cs}{\mathrm{cs}}
\title{The additive registration law for restricted-digit rails}
\author{Samuel Lavery}
\date{Draft --- \today}
\begin{document}
\maketitle

\begin{abstract}
Let $p_1,\dots,p_K$ be distinct odd primes, $d_i\ge1$, and let $X_i$ be the indicator
that the first $d_i$ base-$p_i$ digits of $n$ all lie in
$D_{p_i}=\{0,\dots,\frac{p_i-1}2\}$.  The Chinese remainder theorem gives the joint
count with relative error $O(M/N)$, $M=\prod_ip_i^{d_i}$ --- vacuous once $M>N$.  We
show the true error is governed not by $M$ but by a sum over rails.  The mechanism is
harmonic and exact: \emph{each digit constraint is a square wave}, so each rail's
kernel splits into two lanes on a quarter-turn of phase, $\csc$ on odd frequencies and
$\sec$ on even ones.  The odd lane carries the square wave's own harmonics; the even
lane is nothing but the half-unit registration offset $|D_p|=\frac p2+\frac12$, and is
uniformly of size $\frac1{2p}$ --- a fact no $1/|\alpha|$ envelope can see, and the one
that makes the even lane provably additive with no Diophantine input at all.  We prove
the single-mode term, the even lane, and the main term of the odd lane are each
additive, and isolate what remains.
\end{abstract}

\section{The decomposition}

Fix $p_1,\dots,p_K$ distinct odd primes, $d_i\ge1$, $P_i=p_i^{d_i}$, $M=\prod_iP_i$.
Let $c_i:\mathbb Z/P_i\to\{0,1\}$ be the indicator of the depth-$d_i$ restricted-digit
set and $\rho_i=\hat c_i(0)=\bigl(\frac{p_i+1}{2p_i}\bigr)^{d_i}$ its density.  Write
$S_N(\theta)=\sum_{n\le N}\e{n\theta}$, so $|S_N(\theta)|\le\min(N,\tfrac1{2\|\theta\|})$.

\begin{lemma}[Exact decomposition]\label{lem:dec}
For every $N\ge1$,
\[
\sum_{n\le N}\prod_{i=1}^K X_i(n)\;-\;N\prod_i\rho_i
\;=\;\sum_{\alpha\ne0}\ \prod_{i=1}^K\hat c_i(\alpha_i)\ S_N(\theta_\alpha),
\qquad
\theta_\alpha=\sum_i\frac{\alpha_i}{P_i}=\frac{A(\alpha)}{M},
\]
the sum over $\alpha=(\alpha_1,\dots,\alpha_K)\in\prod_i\mathbb Z/P_i$, where
$A(\alpha)\equiv\sum_i\alpha_i\prod_{j\ne i}P_j \pmod M$.  The map $\alpha\mapsto A$
is a bijection $\prod_i\mathbb Z/P_i\to\mathbb Z/M$, and $A(\alpha)=0$ only for
$\alpha=0$.
\end{lemma}

\begin{proof}
Expand each $c_i$ in characters of $\mathbb Z/P_i$ and exchange the order of summation;
the $\alpha=0$ term is $N\prod_i\rho_i$.  Bijectivity of $\alpha\mapsto A$ is the
Chinese remainder theorem, the $P_i$ being pairwise coprime.
\end{proof}

So $\|\theta_\alpha\|=|A|/M$ is the distance from the mode $\alpha$ to the common mode,
and $|S_N(\theta_\alpha)|\le\min(N,M/2|A|)$ weights each mode by how near it resonates
with the carrier.  Verified: for $(p,d)=(3,3),(5,2)$ and $(3,4),(5,3)$ and
$(7,2),(11,2)$ the right-hand side reproduces the exact error to four decimals.

\section{Harmonization: the digit constraint is a square wave}

Everything below rests on one observation, which is exact.

\begin{lemma}[Common mode plus square wave]\label{lem:sq}
For every odd prime $p$,
\[
c_p(a)\;=\;\tfrac12\bigl(1+\sigma_p(a)\bigr),\qquad
\sigma_p(a)=\begin{cases}+1,&0\le a\le\frac{p-1}2,\\[2pt]-1,&\frac{p+1}2\le a\le p-1.\end{cases}
\]
The constant $\tfrac12$ is carried identically by every rail at every depth: it is the
common mode, and it is why the joint density is $2^{-\sum_id_i}$ up to $O(1/p_i)$.  All
deviation lives in the square waves $\sigma_{p_i}$.
\end{lemma}

A square wave has only odd harmonics, and the entire deviation from that rule is the
half-unit by which $|D_p|=\frac{p+1}2$ exceeds $\frac p2$.  Both statements are visible
at once in closed form.

\begin{lemma}[Two lanes on the quarter-turn]\label{lem:lanes}
Let $|\alpha|\le\frac{p-1}2$, $\alpha\ne0$, and set the phase
$x=\dfrac{\pi\alpha}{2p}\in\bigl(0,\tfrac\pi4\bigr]$.  Then \emph{exactly}
\[
\bigl|\hat c_p(\alpha)\bigr|\;=\;\frac{1}{2p}\cdot
\begin{cases}
\csc x, & \alpha\ \text{odd}\quad\text{(the square wave's own harmonics)},\\[4pt]
\sec x, & \alpha\ \text{even}\quad\text{(the half-unit registration offset)}.
\end{cases}
\]
Consequently
\[
\bigl|\hat c_p(\alpha)\bigr|\le\frac1{2|\alpha|}\ \ (\alpha\ \text{odd}),
\qquad
\bigl|\hat c_p(\alpha)\bigr|\le\frac{1}{\sqrt2\,p}\ \ (\alpha\ \text{even}).
\]
\end{lemma}

\begin{proof}
$\hat c_p(\alpha)=\frac1p\sum_{a<D}\e{-a\alpha/p}$ with $D=\frac{p+1}2$, so
$|\hat c_p(\alpha)|=\frac1p\bigl|\sin(\pi D\alpha/p)/\sin(\pi\alpha/p)\bigr|$.  Now
$\frac{D\alpha}{p}=\frac\alpha2+\frac{\alpha}{2p}=\frac\alpha2+\frac x\pi$, so
$|\sin(\pi D\alpha/p)|$ equals $\cos x$ for $\alpha$ odd and $\sin x$ for $\alpha$ even.
Divide by $|\sin(\pi\alpha/p)|=2\sin x\cos x$.  For the corollaries use
$\sin x\ge\frac{2x}\pi=\frac{|\alpha|}p$ and, since $x\le\frac\pi4$,
$\cos x\ge\frac1{\sqrt2}$.
\end{proof}

The second bound is the payoff.  On even frequencies the kernel has \emph{no} $1/|\alpha|$
decay to give --- it is flat at scale $1/p$ all the way out.  At $\alpha=2$ the
harmonized value is $\frac1{\sqrt2\,p}$ where the envelope $\frac1{2|\alpha|}$ charges
$\frac14$: an overcharge by a factor $\asymp p$.  The even lane exists only because $p$
is odd, i.e.\ because the digit set books a half-unit; charge it as a half-unit and it
costs nothing.

\begin{lemma}[Addition law]\label{lem:add}
For all $x,y$ with $\sin x\sin y\sin(x+y)\ne0$,
\[
\csc x\,\csc y\;=\;\frac{\cot x+\cot y}{\sin(x+y)},
\qquad\text{and}\qquad
\prod_{i=1}^K\csc x_i=\frac{1}{\sin\bigl(\sum_ix_i\bigr)}\prod_{m=2}^{K}
\Bigl(\cot\bigl(x_1+\dots+x_{m-1}\bigr)+\cot x_m\Bigr).
\]
\end{lemma}

\begin{proof}
$\cot x+\cot y=\frac{\cos x\sin y+\cos y\sin x}{\sin x\sin y}=\frac{\sin(x+y)}{\sin x\sin y}$;
the $K$-fold form follows by induction on $m$.  Both verified to machine precision.
\end{proof}

The product of two rails' amplitudes is the \emph{sum} of their cotangents, divided by
the amplitude at the registered phase $x+y$ --- and $x+y$ is exactly the carrier phase
$\theta_\alpha$ at which $S_N$ is read, since $\theta_\alpha=\frac2\pi(x+y)$.  That is
the additive law in one identity.  Its use, however, must be through a reorganisation
of the summation and not term by term; see the register in \S5.

\section{The single-mode term}

Let $\alpha$ have support $\{i\}$.  Summing Lemma~\ref{lem:dec} over such $\alpha$ gives
$\bigl(\prod_{j\ne i}\rho_j\bigr)\cdot E_i$ with
$E_i=\sum_{\alpha\ne0}\hat c_i(\alpha)S_N(\alpha/P_i)$ the single-rail error.

\begin{lemma}[Single-mode bound]\label{lem:single}
For $d=1$, $|E_i|\le\frac{\pi^2}{12}\,p_i+O(1)$.  For $d\ge2$,
$|E_i|\ll p_i(\log p_i)^{d_i-1}$.
\end{lemma}

\begin{proof}
For $d=1$, Lemma~\ref{lem:lanes} gives $|\hat c(\alpha)|\le\frac1{2|\alpha|}$, and
$|S_N(\alpha/p)|\le\frac{p}{2|\alpha|}$, so
$|E|\le2\sum_{1\le\alpha\le p/2}\frac1{2\alpha}\cdot\frac{p}{2\alpha}\le\frac{\pi^2}{12}p$.
For $d\ge2$ factor $\hat c$ into one Dirichlet kernel per digit,
$\hat c(\alpha)=\prod_{j<d}\frac1p\sum_{a\in D_p}\e{-a\alpha/p^{\,d-j}}$; with
$\alpha=\sum_i\alpha_ip^i$ the $j$-th factor is $\ll1/\alpha^{(j)}$ and
$|S_N|\ll p^d/|\alpha|$ with $|\alpha|\asymp\alpha_{d-1}p^{d-1}$.  The top digit
contributes $\sum\alpha_{d-1}^{-2}=O(1)$ and each remaining digit a factor
$\sum_{\alpha_j\le p}\alpha_j^{-1}\ll\log p$.
\end{proof}

Summing over $i$: the single-mode contribution is
$\ll\sum_ip_i(\log p_i)^{d_i-1}\le\sum_iP_i$.  \emph{This is the additive main term.}
As an upper bound it is of the right order and lossy by the cancellation it discards:
with the constant $\pi^2/6$, $(p,d)=(3,9),(5,6)$ predicts $99.3$ against an observed
$27$; $(3,7),(5,5),(7,4)$ predicts $148.7$ against $78$; $(3,7),(5,5),(7,4),(11,3)$
predicts $252.7$ against $4$.

\section{Registration contracts}

The product structure of \S1 is a chart artifact of expanding all rails at once.
Register them \emph{one at a time} instead and the defect never compounds.

\begin{theorem}[Contraction law]\label{thm:contract}
Let $\mathcal E_k=\bigl|\sum_{n\le N}\prod_{i\le k}X_i(n)-N\prod_{i\le k}\rho_i\bigr|$ and
\[
R_k\;=\;\sum_{n\le N}\Bigl(\prod_{i<k}X_i(n)\Bigr)\,\tau_k(n),
\qquad \tau_k:=2\bigl(X_k-\rho_k\bigr),
\]
so that $\tau_k$ is the mean-zero square wave of rail $k$ (Lemma~\ref{lem:sq}).  Then
\[
\mathcal E_k\;\le\;\rho_k\,\mathcal E_{k-1}+\tfrac12|R_k|
\qquad\text{and hence}\qquad
\mathcal E_K\;\le\;\sum_{m\le K}\Bigl(\prod_{m<i\le K}\rho_i\Bigr)\tfrac12|R_m|
\;\le\;\max_{m\le K}|R_m| .
\]
\end{theorem}

\begin{proof}
Substituting $X_k=\rho_k+\frac12\tau_k$ into $\sum_n\prod_{i\le k}X_i$ splits it as
$\rho_k\sum_n\prod_{i<k}X_i+\frac12R_k$; subtract $N\prod_{i\le k}\rho_i
=\rho_k\cdot N\prod_{i<k}\rho_i$ and take absolute values.  Iterating gives the
weighted sum, and $\rho_i\le\frac23$ for $p_i\ge3$ makes the geometric factors sum to
at most $2\cdot\frac12=1$.
\end{proof}

\emph{Each new rail multiplies the accumulated defect by $\rho\approx\frac12$ and adds
one fresh correlation.}  That is the whole additive law: defects register by
contraction, not by multiplication, which is why $K$ rails cost $K$ terms and not
$N^{K}$.  Measured over eight configurations at $N=6\times10^4$ (from two rails to six,
depths $1$ to $5$), $\mathcal E_K/\max_m|R_m|$ lies in $[0.262,0.654]$ --- never above
$1$, and clustering near $\frac12$ as the recursion predicts.

\begin{corollary}[Order the registration]\label{cor:order}
The bound weights rail $m$ by $\prod_{i>m}\rho_i\approx2^{-(K-m)}$.  Since the rails may
be registered in any order, the largest correlations should be registered
\emph{first}, where they are damped hardest.  For $3^5,5^3,7^2,11^2,13$ the ladder
$|R_m|=5.5,\,10.5,\,5.5,\,25.1,\,1.7$ puts its maximum fourth; moving that rail to the
front damps it by $2^{-3}$.
\end{corollary}

This replaces the lattice analysis of the next section entirely.  What remains is a
single object, $R_m$: \emph{one} mean-zero square wave correlated against a fixed
periodic set --- no product over rails, and no Diophantine input.

\section{The cross-mode term for a pair}

For completeness we record what the two-rail Fourier sum says directly, since it fixes
the constant and shows where a pair can be expensive.  Take $K=2$ at depth $1$, $p<q$,
$M=pq$; write $A=\alpha q+\beta p$, so $\|\theta_\alpha\|M=\min(|A|,M-|A|)$, and split
the double sum by the lane of $\alpha$.

\begin{proposition}[The even lane is additive, unconditionally]\label{prop:even}
\[
\Sigma_{\rm ev}:=\sum_{\substack{\alpha\ \text{even}\\ \alpha\ne0}}\sum_{\beta}
\bigl|\hat c_p(\alpha)\bigr|\bigl|\hat c_q(\beta)\bigr|\min\Bigl(N,\frac{M}{2|A|}\Bigr)
\;\ll\;q\log q+\log p\log q .
\]
\end{proposition}

\begin{proof}
Fix $\beta$.  As $\alpha$ runs over the even residues, $A=\beta p+2q\alpha'$ runs over an
arithmetic progression of common difference $2q$ with distinct terms mod $M$; hence the
$m$-th nearest to $0$ has $|A|\ge2(m-1)q$, and
\[
\sum_{\alpha\ \text{even}}\min\Bigl(N,\frac{M}{2|A|}\Bigr)
\;\le\;\min\Bigl(N,\frac M2\Bigr)+2\sum_{1\le m\le p/4}\frac{M}{4mq}
\;\ll\;\min\Bigl(N,\frac M2\Bigr)+p\log p .
\]
Now apply the harmonized even-lane bound $|\hat c_p(\alpha)|\le\frac1{\sqrt2\,p}$ of
Lemma~\ref{lem:lanes}, uniformly in $\alpha$, and $\sum_\beta|\hat c_q(\beta)|\ll\log q$:
\[
\Sigma_{\rm ev}\;\ll\;\frac{\log q}{p}\Bigl(\min\bigl(N,\tfrac{pq}2\bigr)+p\log p\Bigr)
\;\ll\;q\log q+\log p\log q,
\]
since $\frac1p\min(N,\frac{pq}2)\le\frac q2$ in every regime of $M/N$.
\end{proof}

Note what carried the proof: the uniform $\frac1{2p}$ scale of the even lane, with no
Diophantine input whatever.  Under the $\frac1{2|\alpha|}$ envelope the same argument
returns the weighted min-sum and stalls.

On the odd lane the amplitude is genuinely large at small $\alpha$ --- it is the square
wave's fundamental --- and the argument is a lattice count.  By Lemma~\ref{lem:lanes},
\begin{equation}\label{eq:odd}
\Sigma_{\rm odd}\;\le\;\frac14\sum_{\alpha,\beta\ne0}\frac{1}{|\alpha\beta|}
\min\Bigl(N,\frac{M}{2|A|}\Bigr).
\end{equation}
For fixed $\beta$ the minimum of $|A|$ over $\alpha$ is $h(\beta):=|\beta p \bmod^{\pm}q|$,
and $\beta\mapsto h$ is two-to-one onto $\{1,\dots,\frac{q-1}2\}$; on the resonance line
$|\alpha|\asymp|\beta|p/q$.

\begin{proposition}[Odd lane, main term]\label{prop:odd}
The contribution to \eqref{eq:odd} of the modes lying in dyadic boxes
$|\beta|\sim B$, $h\sim H$ with lattice-typical occupancy $\ll BH/q$ is
$\ll q\log^2M$.
\end{proposition}

\begin{proof}
On the resonance line the weight is $\frac1{|\alpha\beta|}\asymp\frac{q}{p\beta^2}$, so
that contribution is $\asymp\frac q{2p}\sum_h\beta_h^{-2}\min(N,\frac M{2h})$ with
$\beta_h=hp^{-1}\bmod^\pm q$.  Summing dyadically with occupancy $BH/q$,
\[
\sum_{B,H}\frac{1}{B^2}\min\Bigl(N,\frac M{2H}\Bigr)\frac{BH}{q}
\;\le\;\frac1q\sum_B\frac1B\sum_H H\min\Bigl(N,\frac M{2H}\Bigr)\;\ll\;\frac{M\log M}{q}=p\log M,
\]
using $H\min(N,\frac M{2H})\le\frac M2$; multiplying by $\frac q{2p}$ gives $\ll q\log M$.
The off-resonance modes ($|A|\ge$ one full spacing) contribute
$\sum_{\alpha}\frac1{|\alpha|}\cdot\frac{p\log q}{|\alpha|}\ll p\log q$ by the same
argument applied to the complementary range.
\end{proof}

\begin{remark}[The beat, and why it costs nothing here]\label{rem:beat}
A pair can be genuinely expensive.  The mode $(\alpha,\beta)=(1,-1)$ sits at
$\|\theta\|=\frac{|q-p|}{M}$, so two rails of nearly equal period \emph{beat}, and
\[
\sup_{N}\Bigl|\sum_{n\le N}X_pX_q-N\rho_p\rho_q\Bigr|
\;\le\;\frac{1.11}{\pi^3}\,\mathcal R(p,q),
\qquad
\mathcal R(p,q)=\max_{\alpha,\beta\ne0}\frac{1}{|\alpha\beta|\,\bigl\|\frac\alpha p+\frac\beta q\bigr\|},
\]
measured over fourteen pairs with the ratio in $[0.13,1.10]$ and equal to
$0.98$--$1.10$ whenever the maximising mode is $(1,-1)$.  For twin primes this is
$\asymp\frac{pq}{2\pi^3}$ and does exceed $C(p+q)$: at $(1481,1483)$ the supremum is
$34\,642=11.7\,(p+q)$.  One $\frac1\pi$ comes from each rail's square-wave
fundamental and one from the carrier resolving the beat.

None of this is needed.  Since $|\alpha\beta|\ge1$ and $\|\theta\|M\ge1$ one has
$\mathcal R(p,q)\le pq$ \emph{unconditionally}, so the pair costs at most its own CRT
modulus --- never the $K$-fold product.  The budget is \textbf{pairwise}: it asks
$P_iP_j\lesssim N$ for each pair, not $\prod_iP_i\lesssim N$.  With rails placed at
$P_i\asymp N^{1/4}$ every pair sits at $N^{1/2}$, half the exponent the budget allows,
and no continued fraction, lattice minimum or prime gap enters anywhere.  The
Diophantine worst case is real but lies outside the configurations \S6 uses.
\end{remark}

\section{Statement, and its register}

\begin{theorem}[Registration law --- \emph{measured}]\label{thm:add}
With the notation above, for all $N\ge1$,
\[
\Bigl|\sum_{n\le N}\prod_iX_i(n)\;-\;N\prod_i\rho_i\Bigr|
\;\le\;C\Bigl(\sum_{i}P_i+\sum_{i<j}\min\bigl(N,P_iP_j\bigr)\Bigr),
\]
with $C\le3.2$ over the tested range.  The bare form $C\sum_iP_i$, which earlier drafts
of this note asserted, is \textbf{false}: two rails of nearly equal period beat, and at
$(p,q)=(1481,1483)$ the supremum over $N$ is $11.7\,(p+q)$ and unbounded in $p+q$
(Remark~\ref{rem:beat}).  The pair term repairs it, and is inert for the rail
placements of \S6.
\end{theorem}

\emph{Proved.}  Lemmas~\ref{lem:dec}, \ref{lem:sq}, \ref{lem:lanes}, \ref{lem:add},
\ref{lem:single}; Theorem~\ref{thm:contract} and Corollary~\ref{cor:order};
Propositions~\ref{prop:even} and~\ref{prop:odd}.  Together these give: the square-wave
decomposition and the two lanes exactly; the additive bound for the single-mode term
and for the entire even lane; the pairwise (never $K$-fold) ceiling $\mathcal R\le P_iP_j$;
and the contraction $\mathcal E_K\le\max_m|R_m|$, which is what removes the product from
the problem.

\emph{Measured.}  Theorem~\ref{thm:add}: $38$ configurations at $N=10^7$,
$M/N\in[10^{-3},2\times10^{14}]$, $|{\rm err}|/\sum_iP_i$ with minimum
$2.6\times10^{-5}$, median $9.5\times10^{-3}$, maximum $3.19$; by contrast
$|{\rm err}|/M$ spans $18$ orders of magnitude and predicts nothing.  Separately, the
Fourier bound \eqref{eq:odd} itself was evaluated exactly for ten prime pairs up to
$(401,409)$ at $N=1.2\times10^4$ with $M/N\in[0.01,13.7]$: the ratio to $p+q$ stayed in
$[0.05,4.24]$.  \emph{The Fourier bound is already additive; no cancellation beyond it
is needed.}

\emph{Rejected, with the measurement that killed each.}
\begin{enumerate}\itemsep2pt
\item The joint error is the sum of single-rail errors --- refuted, the cross modes
carry essentially all of it, $|{\rm cross}|/|{\rm joint}|\approx1$.
\item The defect concentrates at the common mode --- refuted, $|A|=1$ carries $0.0\%$,
$1.3\%$, $0.3\%$ of the mass in the three enumerable cases.
\item $\prod_iL_i\cdot\log M$ suffices, $L_i=\sum_\alpha|\hat c_i(\alpha)|$ --- refuted,
fails by $46\times$ and $95\times$ for large primes at depth~$1$.
\item \textbf{The addition law used term by term.}  Bounding
$\csc x\csc y\,|S_N|\le(|\cot x|+|\cot y|)|S_N|/|\sin(x+y)|$ and summing gives, at
$(101,103)$, $6.9\times10^4$ against the product bound's $581$ --- worse by $119\times$,
and worse on all ten pairs tested.  The identity is exact, but the triangle inequality
discards precisely the reorganisation that makes it additive: the gain lives in
$\sum_\alpha\cot x_\alpha\sum_\beta\Psi(x_\alpha+y_\beta)$, one rail's mass against the
other rail's \emph{complete} phase sum, not in any term-by-term majorant.
\item \textbf{The first minimum $\lambda_1$ of the lattice
$\{(a,r):aP_2\equiv r\ (\mathrm{mod}\ P_1)\}$ controls the exceptional term.}  Refuted
by construction: $109\equiv1\ (\mathrm{mod}\ 27)$, $487\equiv1\ (\mathrm{mod}\ 243)$ and
$251\equiv1\ (\mathrm{mod}\ 125)$ all give $\lambda_1=1$, the shortest possible vector,
yet their measured ratios $|{\rm err}|/(P_1+P_2)$ are $0.072$, $0.012$ and $0.049$ ---
cleaner than the generic mean.  A short vector is not a defect; the bound that said so
was lossy, not the phenomenon.
\end{enumerate}

\emph{Remaining --- one object.}  By Theorem~\ref{thm:contract} everything reduces to
\[
R_m=\sum_{n\le N}\Bigl(\prod_{i<m}X_i(n)\Bigr)\tau_m(n),
\]
a \emph{single} mean-zero square wave correlated against a fixed set periodic mod
$Q=\prod_{i<m}P_i$, with $\gcd(Q,P_m)=1$.  It vanishes identically over any complete
period $QP_m$, so it is a pure incomplete-period defect; and being one correlation
rather than a product over rails, it carries no $K$-fold cost.  Bounding
$\max_m|R_m|$ is the whole remaining step.  Two routes are open and both are
elementary: Fourier on $\tau_m$ alone, giving
$|R_m|\le\sum_{\gamma\ne0}|\hat\tau_m(\gamma)|\,|G_N(\gamma/P_m)|$ with the two-lane
bounds of Lemma~\ref{lem:lanes} on $\hat\tau_m$; or fibring $n$ by its residue mod $Q$,
which turns the inner sum into a mean-zero square wave sampled along an arithmetic
progression of step $Q$ modulo $P_m$.

\begin{remark}[Consequences, pending the proof]
Theorem~\ref{thm:add} replaces the budget $\prod_ip_i^{d_i}\le N$ by
$\sum_ip_i^{d_i}\le N$.  Under the additive budget every rail may be resolved to depth
$(1-\eta)k_i$ regardless of how many rails are used, since $K$ rails cost $K\cdot N^{1-\eta}$
rather than $N^{K(1-\eta)}$; only band~$1$ need be discarded, at Mertens cost $\log2$
instead of $\log K$.  Consequently, in the companion note: the constant $2.36$ becomes
$\approx1.24$; the trade between the bound ($\log K$) and the exceptional set
($\lambda^{-K}$) disappears on the bound side, so $K$ may be taken unbounded; and the
$\log\log\log n$ of the all-$n$ theorem, which rests on the same discard, must be
re-derived.  Those three statements carry a flag until Theorem~\ref{thm:add} is proved.
\end{remark}

\end{document}
